%% ----------------------------------------------------------------------------
%%  §7.7  Appendice: differenziabilità
%% ----------------------------------------------------------------------------
\section{Appendice: differenziabilità}
\label{sec:differenziabilita}

È ben noto che, per le funzioni di una variabile reale, la derivabilità in un punto
implica la continuità della funzione nel punto (\cref{teo:derivabile-continua}). Non vale
una proprietà analoga per le derivate parziali: esistono funzioni $f(x,y)$, definite in un
intorno di un punto $(x_{0},y_{0})$, che hanno le derivate parziali
$f_{x}(x_{0},y_{0})$, $f_{y}(x_{0},y_{0})$ pur non essendo continue in $(x_{0},y_{0})$.
Consideriamo infatti la funzione
\begin{equation}\label{eq:f-discontinua}
  f(x,y)=\frac{xy}{x^{2}+y^{2}} \quad\text{se }(x,y)\neq(0,0), \qquad f(0,0)=0 .
\end{equation}
La funzione si annulla, oltre che in $(0,0)$, in tutti gli altri punti degli assi
coordinati; pertanto
\[
  f_{x}(0,0)=\lim_{h\to0}\frac{f(h,0)-f(0,0)}{h}=\lim_{h\to0}0=0,
\]
e analogamente $f_{y}(0,0)=0$. Però $f$ non è continua in $(0,0)$, perché non ha limite
per $(x,y)\to(0,0)$: lungo la retta $y=x$ vale costantemente $\frac12$, lungo la retta
$y=-x$ vale $-\frac12$ (\cref{es:xy-su-x2y2,fig:discontinua}). Il motivo è chiaro: le
derivate parziali descrivono il comportamento di $f$ soltanto lungo due rette, e non
vedono che cosa accade nelle altre direzioni. La nozione naturale per le funzioni di più
variabili, che estende davvero il concetto di derivabilità, è quella di
\emph{differenziabilità}.

\begin{definizione}[Funzione differenziabile]\label{def:differenziabile}
\index{funzione!differenziabile}\index{differenziale}
Una funzione $f$, definita in un intorno del punto $(x,y)$, è \emph{differenziabile} in
$(x,y)$ se esistono in tale punto le derivate parziali $f_{x}$, $f_{y}$ e se risulta
\begin{equation}\label{eq:differenziabile}
  \lim_{(h,k)\to(0,0)}\frac{f(x+h,y+k)-f(x,y)-f_{x}(x,y)\,h-f_{y}(x,y)\,k}{\sqrt{h^{2}+k^{2}}}=0 .
\end{equation}
In tal caso la funzione lineare $(h,k)\mapsto f_{x}(x,y)\,h+f_{y}(x,y)\,k$ si chiama il
\emph{differenziale} di $f$ in $(x,y)$. Indicando con $(x_{0},y_{0})$ il punto e con
$(x,y)$ il punto variabile, il piano
\begin{equation}\label{eq:piano-tangente}
  z=f(x_{0},y_{0})+f_{x}(x_{0},y_{0})(x-x_{0})+f_{y}(x_{0},y_{0})(y-y_{0})
\end{equation}
si chiama \emph{piano tangente} al grafico di $f$ nel punto $(x_{0},y_{0},f(x_{0},y_{0}))$.
\end{definizione}

La \eqref{eq:differenziabile} è la generalizzazione della
\cref{prop:approssimazione-lineare}: dice che
$f(x+h,y+k)=f(x,y)+f_{x}h+f_{y}k+o\bigl(\sqrt{h^{2}+k^{2}}\bigr)$, cioè che vicino al punto il
grafico di $f$ si confonde con il piano tangente, con un errore trascurabile rispetto alla
distanza. Si può dare una definizione apparentemente più generale, senza richiedere a
priori l'esistenza delle derivate parziali, chiedendo che esistano due numeri $a$, $b$
con $f(x+h,y+k)=f(x,y)+ah+bk+o\bigl(\sqrt{h^{2}+k^{2}}\bigr)$; ponendo $k=0$ si vede subito che
necessariamente $a=f_{x}(x,y)$, e analogamente $b=f_{y}(x,y)$.

\begin{teorema}[La differenziabilità implica la continuità]\label{teo:differenziabile-continua}
Ogni funzione differenziabile in un punto $(x,y)$ è continua in $(x,y)$.
\end{teorema}

\begin{schema}
\Ipotesi{$f$ differenziabile in $(x,y)$.}
\Tesi{$f(x+h,y+k)\to f(x,y)$ per $(h,k)\to(0,0)$.}
\end{schema}

\begin{strategia}
Si procede come per il \cref{teo:derivabile-continua}: l'incremento $f(x+h,y+k)-f(x,y)$ si
scrive come la somma del resto della \eqref{eq:differenziabile}, moltiplicato per
$\sqrt{h^{2}+k^{2}}$, e del differenziale $f_{x}h+f_{y}k$. Entrambi tendono a zero.
\end{strategia}

\begin{dimostrazione}
Poniamo $\rho=\sqrt{h^{2}+k^{2}}$ e, per $(h,k)\neq(0,0)$,
\[
  \varepsilon(h,k)=\frac{f(x+h,y+k)-f(x,y)-f_{x}(x,y)\,h-f_{y}(x,y)\,k}{\rho},
\]
che tende a $0$ per l'ipotesi di differenziabilità. Allora
\[
  \abs{f(x+h,y+k)-f(x,y)}=\Bigl|\varepsilon(h,k)+f_{x}\frac h\rho+f_{y}\frac k\rho\Bigr|\,\rho
  \le\bigl\{\abs{\varepsilon(h,k)}+\abs{f_{x}}+\abs{f_{y}}\bigr\}\,\rho,
\]
perché $\abs h\le\rho$ e $\abs k\le\rho$. Per $(h,k)\to(0,0)$ la quantità tra parentesi
graffe è limitata e $\rho\to0$: il secondo membro tende a zero, e quindi
$f(x+h,y+k)\to f(x,y)$, che è la continuità di $f$ nel punto $(x,y)$.
\end{dimostrazione}

La funzione \eqref{eq:f-discontinua} non è continua nell'origine, e quindi, per il
teorema appena dimostrato, non vi è differenziabile, benché abbia le derivate parziali.

\subsection{Derivate direzionali e significato del gradiente}

\begin{teorema}[Derivate direzionali delle funzioni differenziabili]\label{teo:derivata-direzionale}
\index{derivata!direzionale}
Se $f$ è differenziabile nel punto $P=(x,y)$, allora per ogni versore
$v=(\cos\vartheta,\sen\vartheta)$ esiste la derivata direzionale \eqref{eq:derivata-direzionale},
e risulta
\begin{equation}\label{eq:direzionale-gradiente}
  \frac{\partial f}{\partial v}(P)=f_{x}(P)\cos\vartheta+f_{y}(P)\sen\vartheta .
\end{equation}
Di conseguenza, se $\nabla f(P)\neq(0,0)$, la derivata direzionale è massima, e uguale a
$\abs{\nabla f(P)}$, nella direzione del gradiente, ed è nulla nelle direzioni
perpendicolari al gradiente.
\end{teorema}

\begin{schema}
\Ipotesi{$f$ differenziabile in $P$; $v=(\cos\vartheta,\sen\vartheta)$.}
\Tesi{$\frac{\partial f}{\partial v}(P)=\nabla f(P)\cdot v$; massima nella direzione di
$\nabla f(P)$.}
\end{schema}

\begin{strategia}
Basta scegliere, nella \eqref{eq:differenziabile}, l'incremento $(h,k)=tv$, che ha
lunghezza $\abs t$: il rapporto incrementale lungo la direzione $v$ differisce da
$f_{x}\cos\vartheta+f_{y}\sen\vartheta$ per una quantità infinitesima. La seconda parte è poi
la stessa trigonometria della \cref{prop:gradiente-lineare}.
\end{strategia}

\begin{dimostrazione}
Con $(h,k)=(t\cos\vartheta,t\sen\vartheta)$, $t\neq0$, si ha $\sqrt{h^{2}+k^{2}}=\abs t$ e
\begin{multline*}
  \frac{f(x+t\cos\vartheta,y+t\sen\vartheta)-f(x,y)}{t}-\bigl(f_{x}\cos\vartheta+f_{y}\sen\vartheta\bigr)\\
  =\frac{\abs t}{t}\,\varepsilon(t\cos\vartheta,t\sen\vartheta),
\end{multline*}
dove $\varepsilon$ è la funzione della dimostrazione precedente. Il punto
$(t\cos\vartheta,t\sen\vartheta)$ dista $\abs t$ dall'origine; quindi, per la definizione di
limite \eqref{eq:limite-piano-simboli} applicata alla \eqref{eq:differenziabile}, per ogni
$\eta>0$ esiste $\delta>0$ tale che $\abs{\varepsilon(t\cos\vartheta,t\sen\vartheta)}<\eta$ non appena
$0<\abs t<\delta$, cioè $\varepsilon(t\cos\vartheta,t\sen\vartheta)\to0$ per $t\to0$; mentre $\frac{\abs t}t=\pm1$ è limitata: il secondo membro tende a zero, e si ottiene la
\eqref{eq:direzionale-gradiente}. Per la seconda parte, scritto
$\nabla f(P)=r(\cos\varphi,\sen\varphi)$ con $r=\abs{\nabla f(P)}>0$, si ha come nella
\cref{prop:gradiente-lineare} $\frac{\partial f}{\partial v}=r\cos(\vartheta-\varphi)$, che è
massima per $\vartheta=\varphi$ e nulla per $\vartheta=\varphi\pm\frac\pi2$.
\end{dimostrazione}

Le direzioni perpendicolari al gradiente sono quelle in cui $f$, al primo ordine, non
varia: sono le direzioni delle tangenti alle curve di livello. Il gradiente è quindi
perpendicolare alle curve di livello e punta verso i valori crescenti di $f$, come si
vede nella \cref{fig:livello} (dove sono disegnati i vettori $-\nabla p$, che puntano
invece verso le pressioni decrescenti). Per la funzione \eqref{eq:f-discontinua}, che non è
differenziabile nell'origine, la formula \eqref{eq:direzionale-gradiente} è falsa: si
avrebbe $\frac{\partial f}{\partial v}(0,0)=0$ in ogni direzione, mentre, ad esempio, lungo
la bisettrice il rapporto incrementale $\frac{f(t/\sqrt2,t/\sqrt2)}{t}=\frac1{2t}$ non ha
limite.

\begin{attenzione}[Le derivate direzionali non bastano]
Il \cref{teo:derivata-direzionale} non si inverte: una funzione può avere in un punto la
derivata direzionale secondo \emph{ogni} direzione senza essere differenziabile, e
nemmeno continua. Riprendiamo la funzione $g(x,y)=\frac{x^{2}y}{x^{4}+y^{2}}$
dell'avvertenza «Le rette non bastano», ponendo $g(0,0)=0$. Per
$v=(\cos\vartheta,\sen\vartheta)$ con $\sen\vartheta\neq0$ si ha
\[
  \frac{g(t\cos\vartheta,t\sen\vartheta)}{t}
  =\frac{\cos^{2}\vartheta\,\sen\vartheta}{t^{2}\cos^{4}\vartheta+\sen^{2}\vartheta},
\]
che tende a $\frac{\cos^{2}\vartheta}{\sen\vartheta}$ per $t\to0$; per $\sen\vartheta=0$
la funzione è nulla e la derivata direzionale vale $0$. Tutte le derivate direzionali
esistono, ma $g$ non è continua nell'origine (vale $\frac12$ lungo la parabola $y=x^{2}$),
e quindi non è differenziabile; coerentemente, la \eqref{eq:direzionale-gradiente} è
falsa, perché $\nabla g(0,0)=(0,0)$. Non basta neppure che la
\eqref{eq:direzionale-gradiente} sia verificata: per la funzione
$\frac{x^{4}y}{x^{8}+y^{2}}$, posta uguale a $0$ nell'origine, tutte le derivate
direzionali nell'origine sono nulle, come il gradiente, eppure lungo la curva $y=x^{4}$
la funzione vale costantemente $\frac12$. Per stabilire la differenziabilità serve
dunque la definizione, oppure il teorema del differenziale che segue.
\end{attenzione}

\subsection{Il teorema del differenziale}

Verificare la differenziabilità con la definizione è spesso laborioso. Un importante
criterio è espresso dal teorema seguente.

\begin{teorema}[del differenziale]\label{teo:differenziale}\index{teorema!del differenziale}
Se $f$ ammette derivate parziali prime in un intorno del punto $(x,y)$, e queste sono
continue in $(x,y)$, allora $f$ è differenziabile in $(x,y)$.
\end{teorema}

\begin{schema}
\Ipotesi{$f_{x}$, $f_{y}$ esistono in un intorno di $(x,y)$ e sono continue in $(x,y)$.}
\Tesi{vale la \eqref{eq:differenziabile}.}
\end{schema}

\begin{strategia}
Andiamo dal punto $(x,y)$ al punto $(x+h,y+k)$ in due tappe, parallele agli assi:
prima da $(x,y)$ a $(x,y+k)$ in verticale, poi da $(x,y+k)$ a $(x+h,y+k)$ in
orizzontale. Su ciascuna tappa la funzione dipende da una sola variabile, e il teorema di
Lagrange esprime l'incremento con una derivata parziale calcolata in un punto
intermedio. La continuità delle derivate parziali garantisce che queste differiscano di
poco da quelle nel punto $(x,y)$.
\end{strategia}

\begin{dimostrazione}
Sia $I_{\delta}$ un intorno di $(x,y)$ in cui esistono $f_{x}$ e $f_{y}$, e sia
$\sqrt{h^{2}+k^{2}}<\delta$. Applicando due volte il teorema di Lagrange (relativo a funzioni
di una variabile) abbiamo
\begin{equation}\label{eq:differenziale-lagrange}
\begin{split}
  f(x+h,y+k)-f(x,y)&=\bigl[f(x+h,y+k)-f(x,y+k)\bigr]\\
  &\quad+\bigl[f(x,y+k)-f(x,y)\bigr]\\
  &=f_{x}(x(h),y+k)\,h+f_{y}(x,y(k))\,k,
\end{split}
\end{equation}
dove $x(h)$ è un punto opportuno dell'intervallo di estremi $x$, $x+h$, mentre $y(k)$ è un
punto opportuno dell'intervallo di estremi $y$, $y+k$ (i segmenti percorsi stanno in
$I_{\delta}$, per lo stesso argomento usato nella dimostrazione del teorema di
Schwarz). Si noti in particolare che $x(h)\to x$ e $y(k)\to y$ per $(h,k)\to(0,0)$. Si
ottiene la stima
\begin{equation}\label{eq:differenziale-stima}
\begin{split}
  &\Bigl|\frac{f(x+h,y+k)-f(x,y)-f_{x}(x,y)h-f_{y}(x,y)k}{\sqrt{h^{2}+k^{2}}}\Bigr|\\
  &\quad\le\abs{f_{x}(x(h),y+k)-f_{x}(x,y)}\frac{\abs h}{\sqrt{h^{2}+k^{2}}}\\
  &\qquad+\abs{f_{y}(x,y(k))-f_{y}(x,y)}\frac{\abs k}{\sqrt{h^{2}+k^{2}}}\\
  &\quad\le\abs{f_{x}(x(h),y+k)-f_{x}(x,y)}+\abs{f_{y}(x,y(k))-f_{y}(x,y)} .
\end{split}
\end{equation}
Per l'ipotesi di continuità delle derivate parziali $f_{x}$ e $f_{y}$ nel punto $(x,y)$,
l'ultimo membro della \eqref{eq:differenziale-stima} tende a zero per $(h,k)\to(0,0)$,
perché i punti $(x(h),y+k)$ e $(x,y(k))$ tendono a $(x,y)$; perciò anche il primo membro
tende a zero.
\end{dimostrazione}

Se una funzione $f$ è continua in un insieme aperto $D$ si dice anche che è \emph{di
classe $C^{0}$ in $D$} e si scrive $f\in C^{0}(D)$; se ammette derivate parziali prime
continue in $D$ si dice che è \emph{di classe $C^{1}$ in $D$} e si scrive $f\in C^{1}(D)$;
più in generale $f\in C^{k}(D)$, con $k\in\N$, significa che $f$ ammette derivate
parziali continue in $D$ fino all'ordine $k$. Infine, se $f\in C^{k}(D)$ per ogni $k$, si
scrive $f\in C^{\infty}(D)$. Con queste notazioni, il teorema del differenziale e la
proprietà di continuità si riassumono così:
\begin{equation}\label{eq:classi-differenziabilita}
  f\in C^{1}(D)\ \implica\ f \text{ differenziabile in } D\ \implica\ f\in C^{0}(D).
\end{equation}
Nessuna delle due implicazioni si inverte. L'\cref{es:non-differenziabile} presenta una
funzione continua non differenziabile. Una funzione differenziabile ovunque ma non di
classe $C^{1}$ è invece $f(x,y)=x^{2}\sen\frac1x$ per $x\neq0$, $f(0,y)=0$: è la funzione $f_{2}$
dell'\cref{oss:f2-non-C1}, pensata come funzione di due variabili. Nei punti $(0,y)$ si ha
$f(h,y+k)-f(0,y)=h^{2}\sen\frac1h$ (o $0$ se $h=0$), che in valore assoluto non supera
$h^{2}\le h^{2}+k^{2}$: con $f_{x}(0,y)=f_{y}(0,y)=0$ la \eqref{eq:differenziabile} è verificata, e
$f$, che fuori dall'asse $y$ è di classe $C^{1}$, è differenziabile ovunque; ma $f_{x}=2x\sen\frac1x-\cos\frac1x$ non è continua nei
punti dell'asse $y$, e quindi $f\notin C^{1}(\R^{2})$.

\begin{esempio}[Propagazione degli errori di misura]\label{es:propagazione-errori}
Il differenziale fornisce la regola pratica per stimare l'errore su una grandezza
calcolata da altre misurate. La densità dell'aria si ricava dalla pressione $p$ e dalla
temperatura assoluta $T$ con l'equazione di stato, $\rho=\frac p{RT}$. Le derivate
parziali sono $\rho_{p}=\frac1{RT}=\frac\rho p$ e $\rho_{T}=-\frac p{RT^{2}}=-\frac\rho T$, continue
per $T>0$; per il teorema del differenziale
\[
  \Delta\rho\approx\rho_{p}\Delta p+\rho_{T}\Delta T, \qquad
  \frac{\Delta\rho}{\rho}\approx\frac{\Delta p}{p}-\frac{\Delta T}{T}, \qquad
  \Bigl|\frac{\Delta\rho}{\rho}\Bigr|\lesssim\frac{\abs{\Delta p}}{p}+\frac{\abs{\Delta T}}{T} :
\]
gli errori relativi, al primo ordine, si sommano. Il programma confronta questa stima
con una simulazione in cui $p$ e $T$ variano a caso nei rispettivi intervalli di
incertezza.
\end{esempio}

\begin{python}[Propagazione degli errori: differenziale e simulazione]
import numpy as np

R = 287.05
rho = lambda p, T: p / (R * T)           # densità dell'aria
p0, T0, dp, dT = 101325.0, 288.15, 300.0, 1.5   # misure e incertezze
r0 = rho(p0, T0)
lin = r0 * (dp / p0 + dT / T0)           # stima con il differenziale
rng = np.random.default_rng(1)
p = p0 + rng.uniform(-dp, dp, 100000)
T = T0 + rng.uniform(-dT, dT, 100000)
print(f"rho = {r0:.5f} kg/m^3")
print(f"errore massimo (differenziale): {lin:.5f}"
      f"  ({lin / r0:.2%})")
print(f"errore massimo osservato (simulazione): "
      f"{np.max(np.abs(rho(p, T) - r0)):.5f}")
\end{python}
\begin{risultato}
rho = 1.22501 kg/m^3
errore massimo (differenziale): 0.01000  (0.82%)
errore massimo osservato (simulazione): 0.01003
\end{risultato}

La stima lineare coincide con l'errore massimo osservato a meno di qualche unità del
quinto decimale: la piccola differenza è dovuta ai termini del second'ordine che il
differenziale trascura (il caso peggiore esatto, con $p=p_{0}+\Delta p$ e $T=T_{0}-\Delta T$,
dà $0{,}010056$).

\subsection{La derivata lungo un segmento e la formula di Taylor}

Per dimostrare il criterio dell'hessiano studiamo $f$ lungo un segmento uscente dal
punto critico: ci servono la derivata prima e seconda di una funzione di due variabili
lungo un segmento, e una formula di Taylor di ordine due con un resto più preciso di
quello di Peano.

\begin{lemma}[Derivate lungo un segmento]\label{lem:derivata-segmento}
Sia $f\in C^{1}\bigl(I_{\delta}(P_{0})\bigr)$, con $P_{0}=(x_{0},y_{0})$, e sia $(h,k)$ con
$0<\sqrt{h^{2}+k^{2}}<\delta$. La funzione di una variabile
\begin{equation}\label{eq:g-segmento}
  g(t)=f(x_{0}+th,\,y_{0}+tk)
\end{equation}
è derivabile in un intervallo aperto contenente $[0,1]$, con
\begin{equation}\label{eq:g-primo}
  g'(t)=f_{x}(P_{t})\,h+f_{y}(P_{t})\,k, \qquad P_{t}=(x_{0}+th,\,y_{0}+tk).
\end{equation}
Se inoltre $f\in C^{2}\bigl(I_{\delta}(P_{0})\bigr)$, allora $g$ è derivabile due volte, $g''$ è
continua e
\begin{equation}\label{eq:g-secondo}
  g''(t)=f_{xx}(P_{t})\,h^{2}+2f_{xy}(P_{t})\,hk+f_{yy}(P_{t})\,k^{2} .
\end{equation}
\end{lemma}

\begin{schema}
\Ipotesi{$f\in C^{1}$ (risp.\ $C^{2}$) in $I_{\delta}(P_{0})$; $0<\sqrt{h^{2}+k^{2}}<\delta$.}
\Tesi{$g'(t)=f_{x}(P_{t})h+f_{y}(P_{t})k$; \ (se $f\in C^{2}$)
$g''(t)=f_{xx}h^{2}+2f_{xy}hk+f_{yy}k^{2}$ in $P_{t}$, continua.}
\end{schema}

\begin{strategia}
La derivata di $g$ è la derivata di $f$ nella direzione del segmento, moltiplicata per la
lunghezza di $(h,k)$: per le funzioni differenziabili (e le funzioni di classe $C^{1}$ lo
sono, per il teorema del differenziale) essa si calcola con il gradiente. La derivata
seconda si ottiene applicando lo stesso risultato alla funzione $hf_{x}+kf_{y}$, e il
teorema di Schwarz riunisce i due termini misti.
\end{strategia}

\begin{dimostrazione}
Posto $r=\sqrt{h^{2}+k^{2}}$, i punti $P_{t}$ stanno in $I_{\delta}(P_{0})$ per
$\abs t<\frac\delta r$, e $\frac\delta r>1$. Sia $v=\frac1r(h,k)$, che è un versore. Per il
teorema del differenziale $f$ è differenziabile in ogni punto $P_{t}$, e per il
\cref{teo:derivata-direzionale}
\begin{align*}
  g'(t)&=\lim_{s\to0}\frac{f(P_{t}+s(h,k))-f(P_{t})}{s}
  =r\lim_{s\to0}\frac{f(P_{t}+sr\,v)-f(P_{t})}{sr}\\
  &=r\,\frac{\partial f}{\partial v}(P_{t})=f_{x}(P_{t})h+f_{y}(P_{t})k .
\end{align*}
Se $f\in C^{2}$, la funzione $\varphi=hf_{x}+kf_{y}$ è di classe $C^{1}$, e $g'(t)=\varphi(P_{t})$.
Applicando a $\varphi$ quanto appena dimostrato,
$g''(t)=h\,\varphi_{x}(P_{t})+k\,\varphi_{y}(P_{t})=h(hf_{xx}+kf_{yx})+k(hf_{xy}+kf_{yy})$, calcolato in
$P_{t}$, che per il teorema di Schwarz ($f_{xy}=f_{yx}$) è la \eqref{eq:g-secondo}. Infine
$g''$ è continua: se $t_{n}\to t$, allora $P_{t_{n}}\to P_{t}$ (le coordinate convergono,
\eqref{eq:distanza-componenti}), e per la continuità delle derivate seconde
(\cref{oss:limite-piano-successioni}) $f_{xx}(P_{t_{n}})\to f_{xx}(P_{t})$, e così per le
altre; quindi $g''(t_{n})\to g''(t)$.
\end{dimostrazione}

\begin{lemma}[Formula di Taylor con il resto di Lagrange]\label{lem:taylor-lagrange}
\index{formula di Taylor!con il resto di Lagrange}
Sia $g$ derivabile due volte in un intervallo aperto contenente $[0,1]$. Esiste
$\tau\in(0,1)$ tale che
\begin{equation}\label{eq:taylor-lagrange}
  g(1)=g(0)+g'(0)+\frac12g''(\tau) .
\end{equation}
\end{lemma}

\begin{schema}
\Ipotesi{$g$ derivabile due volte in un intervallo aperto contenente $[0,1]$.}
\Tesi{esiste $\tau\in(0,1)$ con $g(1)=g(0)+g'(0)+\frac12g''(\tau)$.}
\end{schema}

\begin{strategia}
È la formula di Taylor di ordine $1$ con il resto scritto in un punto intermedio, come
nella formula degli accrescimenti finiti \eqref{eq:accrescimenti-finiti}, che ne è il
caso di ordine $0$. Il resto $g(1)-g(0)-g'(0)$ si confronta con $t^{2}$ mediante il teorema
di Cauchy, che produce un rapporto tra $g'$ e $2t$; un'ulteriore applicazione del
teorema di Lagrange a $g'$ fa comparire la derivata seconda.
\end{strategia}

\begin{dimostrazione}
Applichiamo il teorema di Cauchy (\cref{teo:cauchy}) nell'intervallo $[0,1]$ alle
funzioni $G(t)=g(t)-g(0)-g'(0)\,t$ e $K(t)=t^{2}$, osservando che $K'(t)=2t\neq0$ per
$t\in(0,1)$. Poiché $G(0)=K(0)=0$ e $K(1)=1$, esiste $\tau_{1}\in(0,1)$ tale che
\[
  g(1)-g(0)-g'(0)=\frac{G(1)-G(0)}{K(1)-K(0)}=\frac{G'(\tau_{1})}{K'(\tau_{1})}
  =\frac{g'(\tau_{1})-g'(0)}{2\tau_{1}} .
\]
Per il teorema di Lagrange applicato a $g'$ nell'intervallo $[0,\tau_{1}]$, esiste
$\tau\in(0,\tau_{1})$ con $g'(\tau_{1})-g'(0)=g''(\tau)\,\tau_{1}$; sostituendo si ottiene la
\eqref{eq:taylor-lagrange}.
\end{dimostrazione}

\subsection{La dimostrazione del criterio dell'hessiano}
\label{sec:dim-hessiano}

Siamo ora in grado di dimostrare il \cref{teo:hessiano}. Indichiamo, per ogni punto $P$
vicino a $(x_{0},y_{0})$ e ogni $(h,k)$, la forma quadratica
\begin{equation}\label{eq:forma-quadratica}
  Q_{P}(h,k)=f_{xx}(P)\,h^{2}+2f_{xy}(P)\,hk+f_{yy}(P)\,k^{2} ,
\end{equation}
che per il \cref{lem:derivata-segmento} è la derivata seconda di $f$ lungo il segmento di
direzione $(h,k)$.

\begin{dimostrazione}[Dimostrazione del \cref{teo:hessiano}]
\emph{Passo 1: segno della forma quadratica.} Siano $A=f_{xx}(P)$, $B=f_{xy}(P)$,
$C=f_{yy}(P)$. Se $A\neq0$, completando il quadrato si ha l'identità
\begin{equation}\label{eq:completamento-quadrato}
\begin{split}
  A\,Q_{P}(h,k)&=A^{2}h^{2}+2ABhk+ACk^{2}\\
  &=(Ah+Bk)^{2}+(AC-B^{2})k^{2} .
\end{split}
\end{equation}
Se $A>0$ e $AC-B^{2}=Hf(P)>0$, il secondo membro è positivo per $(h,k)\neq(0,0)$ (se $k=0$
il primo addendo è $A^{2}h^{2}>0$), e quindi $Q_{P}(h,k)>0$.

\emph{Passo 2: il minimo.} Supponiamo $Hf(x_{0},y_{0})>0$ e $f_{xx}(x_{0},y_{0})>0$. Le
funzioni $Hf$ e $f_{xx}$ sono continue (\cref{teo:continue-piano}); per il teorema della
permanenza del segno, che vale per le funzioni continue di due variabili con la stessa
dimostrazione del caso di una variabile, esiste $\delta>0$ tale che $Hf>0$ e $f_{xx}>0$ in
tutto $I_{\delta}(x_{0},y_{0})$, intorno nel quale supponiamo $f$ di classe $C^{2}$. Sia
$0<\sqrt{h^{2}+k^{2}}<\delta$ e sia $g(t)=f(x_{0}+th,y_{0}+tk)$. Per il
\cref{lem:derivata-segmento}, $g'(0)=f_{x}(x_{0},y_{0})h+f_{y}(x_{0},y_{0})k=0$, perché il
punto è critico; per il \cref{lem:taylor-lagrange} esiste $\tau\in(0,1)$ con
\[
  f(x_{0}+h,y_{0}+k)-f(x_{0},y_{0})=g(1)-g(0)=\frac12g''(\tau)=\frac12Q_{P_{\tau}}(h,k)>0,
\]
dove l'ultima disuguaglianza segue dal passo 1, perché $P_{\tau}\in I_{\delta}$. Dunque
$f(x,y)>f(x_{0},y_{0})$ per ogni $(x,y)\neq(x_{0},y_{0})$ in $I_{\delta}$: il punto è di minimo
relativo stretto.

\emph{Passo 3: il massimo.} Se $Hf(x_{0},y_{0})>0$ e $f_{xx}(x_{0},y_{0})<0$, si applica il
passo 2 alla funzione $-f$, che ha lo stesso hessiano (il prodotto $f_{xx}f_{yy}$ e il
quadrato $f_{xy}^{2}$ non cambiano cambiando segno a $f$) e derivata seconda
$(-f)_{xx}=-f_{xx}>0$.

\emph{Passo 4: la sella.} Sia $Hf(x_{0},y_{0})=AC-B^{2}<0$, con $A$, $B$, $C$ calcolati in
$(x_{0},y_{0})$. Troviamo due direzioni lungo le quali la forma quadratica ha segni
opposti. Se $A\neq0$, si prenda $(h,k)=(1,0)$, che dà $Q=A$, e $(h,k)=(-B,A)$, che per la
\eqref{eq:completamento-quadrato} dà $A\,Q=(AC-B^{2})A^{2}$, cioè $Q=A(AC-B^{2})$, di segno
opposto ad $A$. Se $A=0$, allora $-B^{2}<0$, cioè $B\neq0$, e $Q(1,k)=k(2B+Ck)$ ha, per $k$
piccolo, il segno di $2Bk$, che cambia con il segno di $k$. Poiché
$Q(\lambda h,\lambda k)=\lambda^{2}Q(h,k)$, moltiplicando $(h,k)$ per un $\lambda>0$ piccolo il
segno della forma non cambia: possiamo quindi supporre $\sqrt{h^{2}+k^{2}}<\delta$, dove
$I_{\delta}$ è un intorno in cui $f$ ha derivate prime e seconde continue. Sia dunque $(h,k)$
una tale direzione con $Q_{(x_{0},y_{0})}(h,k)>0$, e sia $g(t)=f(x_{0}+th,y_{0}+tk)$: si ha $g'(0)=0$ e
$g''(0)=Q_{(x_{0},y_{0})}(h,k)>0$, con $g''$ continua, e per il criterio della derivata
seconda (\cref{teo:test-derivata-seconda}) $t=0$ è un punto di minimo relativo stretto
di $g$: lungo questa direzione, arbitrariamente vicino a $(x_{0},y_{0})$, la funzione assume
valori maggiori di $f(x_{0},y_{0})$. Lungo una direzione con $Q<0$, allo stesso modo,
assume valori minori. Dunque $(x_{0},y_{0})$ non è né di massimo né di minimo relativo.
\end{dimostrazione}

\begin{notastorica}[Il differenziale da Leibniz a Fréchet]
Per Leibniz il differenziale $dy$ era un incremento «infinitamente piccolo», e la formula
$dz=\frac{\partial z}{\partial x}dx+\frac{\partial z}{\partial y}dy$ fu usata liberamente per
tutto il Settecento. La distinzione tra derivate parziali e differenziabilità non poteva
allora neppure emergere: le funzioni studiate erano espressioni analitiche «regolari», per
le quali le due nozioni vanno di pari passo. Quando, nell'Ottocento, con la revisione
critica avviata da Cauchy, gli infinitesimi furono sostituiti dai limiti e si cominciò a
considerare funzioni qualsiasi, ci si accorse che l'esistenza delle derivate parziali non
basta a giustificare quella formula, e che occorre una condizione più forte. La definizione moderna di
differenziabilità, che chiede che l'incremento sia approssimato da una funzione lineare
con un errore trascurabile rispetto alla distanza, fu formulata da Otto Stolz nel 1893 e
ripresa all'inizio del Novecento da William Henry Young e da Maurice Fréchet, che la
estese agli spazi di funzioni: è la nozione che sta alla base dell'analisi moderna e dei
metodi numerici, dai metodi di Newton in più variabili ai codici di fluidodinamica
computazionale.
\end{notastorica}

\subsection{Approfondimento: funzioni omogenee e teorema di Eulero}
\label{subsec:funzioni-omogenee}
\index{funzione!omogenea}\index{teorema!di Eulero sulle funzioni omogenee}

Molte formule dell'ingegneria sono prodotti di potenze delle variabili: la portanza
$\frac12\rho SV^{2}C_{L}$, o il coefficiente di resistenza indotta
$C_{Di}=\frac{C_{L}^{2}}{\pi e\,AR}$ dell'\cref{es:sensibilita-AR}. Moltiplicando tutte le
variabili per uno stesso fattore $t>0$, la funzione risulta moltiplicata per una potenza
di $t$. Queste funzioni si dicono \emph{omogenee}, e il loro gradiente soddisfa
un'identità notevole, scoperta da Eulero (si veda la nota storica del \S\ref{sec:schwarz}).
Come nel resto del paragrafo, ci limitiamo alle funzioni di due variabili.

\begin{definizione}[Funzione omogenea]\label{def:funzione-omogenea}
Un insieme $D\subseteq\R^{2}\setminus\{(0,0)\}$ si dice un \emph{cono} se, per ogni
$(x,y)\in D$ e ogni $t>0$, anche $(tx,ty)\in D$. Sia $D$ un cono aperto e
sia $\alpha\in\R$. Una funzione $f:D\to\R$ si dice \emph{omogenea di grado $\alpha$} se
\begin{equation}\label{eq:omogenea}
  f(tx,ty)=t^{\alpha}f(x,y) \qquad\perogni(x,y)\in D,\ \perogni t>0 .
\end{equation}
\end{definizione}

Sono omogenei di grado $n$, nel cono $\R^{2}\setminus\{(0,0)\}$, i polinomi i cui monomi
hanno tutti grado $n$, come $x^{2}+3xy$ ($n=2$). La norma $\sqrt{x^{2}+y^{2}}$ e la funzione
$\sqrt{\abs{xy}}$ dell'\cref{es:non-differenziabile} sono omogenee di grado $1$: i loro
grafici sono coni con il vertice nell'origine. La funzione $\frac{xy}{x^{2}+y^{2}}$
dell'\cref{es:xy-su-x2y2} e l'angolo polare $\operatorname{arctg}\frac yx$, nel semipiano
$x>0$, sono omogenee di grado $0$; e una funzione omogenea di grado $0$ è costante su ogni
semiretta uscente dall'origine, perché $f(tx,ty)=f(x,y)$. Il potenziale della doppietta
$\frac{x}{x^{2}+y^{2}}$ ha grado $-1$; il coefficiente $C_{Di}(C_{L},AR)$, definito nel
semipiano $AR>0$, ha grado $2-1=1$.

\begin{teorema}[di Eulero sulle funzioni omogenee]\label{teo:eulero-omogenee}
Sia $D$ un cono aperto e sia $\alpha\in\R$.
\begin{enumerate}[label=\textup{(\roman*)}]
  \item Se $f:D\to\R$ è omogenea di grado $\alpha$ ed è differenziabile nel punto
  $(x,y)\in D$, allora
  \begin{equation}\label{eq:eulero-omogenee}
    x\,f_{x}(x,y)+y\,f_{y}(x,y)=\alpha\,f(x,y) .
  \end{equation}
  \item Viceversa, se $f\in C^{1}(D)$ e la \eqref{eq:eulero-omogenee} vale in ogni punto di
  $D$, allora $f$ è omogenea di grado $\alpha$.
\end{enumerate}
\end{teorema}

\begin{schema}
\Ipotesi{$D$ cono aperto; \ (i) $f$ omogenea di grado $\alpha$, differenziabile in
$(x,y)$; \ (ii) $f\in C^{1}(D)$ e $xf_{x}+yf_{y}=\alpha f$ in tutto $D$.}
\Tesi{(i) $xf_{x}+yf_{y}=\alpha f$ nel punto $(x,y)$; \ (ii) $f(tx,ty)=t^{\alpha}f(x,y)$ per
ogni $(x,y)\in D$ e ogni $t>0$.}
\end{schema}

\begin{strategia}
Tutto si svolge sulla semiretta che esce dall'origine e passa per $(x,y)$, lungo la quale
si studia la funzione di una variabile $\psi(t)=f(tx,ty)$. La sua derivata è, a meno del
fattore $\sqrt{x^{2}+y^{2}}$, la derivata di $f$ nella direzione radiale, che per le
funzioni differenziabili si calcola con il gradiente (\cref{teo:derivata-direzionale}):
si trova $\psi'(t)=xf_{x}(tx,ty)+yf_{y}(tx,ty)$. L'omogeneità dice invece che
$\psi(t)=t^{\alpha}\psi(1)$. Per (i) basta confrontare le due espressioni di $\psi'(1)$; per
(ii) si mostra che $t^{-\alpha}\psi(t)$ ha derivata nulla, e quindi è costante.
\end{strategia}

\begin{dimostrazione}
Sia $(x,y)\in D$, e siano $r=\sqrt{x^{2}+y^{2}}>0$ e
$v=\frac1r(x,y)$, che è un versore. Per $t>0$ il punto $P_{t}=(tx,ty)$ sta in $D$; poniamo
$\psi(t)=f(P_{t})$. Per $s\neq0$ con $t+s>0$ si ha
\[
  \frac{\psi(t+s)-\psi(t)}{s}=\frac{f(P_{t}+s(x,y))-f(P_{t})}{s}
  =r\,\frac{f(P_{t}+sr\,v)-f(P_{t})}{sr} .
\]
Se $f$ è differenziabile in $P_{t}$, per il \cref{teo:derivata-direzionale} l'ultima
frazione tende, per $s\to0$, alla derivata direzionale
$\frac{\partial f}{\partial v}(P_{t})=f_{x}(P_{t})\frac xr+f_{y}(P_{t})\frac yr$; quindi
$\psi$ è derivabile in $t$ e
\begin{equation}\label{eq:psi-radiale}
  \psi'(t)=x\,f_{x}(tx,ty)+y\,f_{y}(tx,ty) .
\end{equation}

(i) Per l'omogeneità $\psi(t)=t^{\alpha}f(x,y)$, da cui $\psi'(1)=\alpha f(x,y)$. D'altra
parte $f$ è differenziabile in $P_{1}=(x,y)$, e la \eqref{eq:psi-radiale} con $t=1$ dà
$\psi'(1)=xf_{x}(x,y)+yf_{y}(x,y)$: confrontando, si ottiene la tesi.

(ii) Per il teorema del differenziale (\cref{teo:differenziale}) $f$ è differenziabile in
ogni punto di $D$, e la \eqref{eq:psi-radiale} vale per ogni $t>0$. Moltiplicandola per
$t$ e applicando l'ipotesi nel punto $P_{t}\in D$, si ottiene
\[
  t\,\psi'(t)=(tx)\,f_{x}(P_{t})+(ty)\,f_{y}(P_{t})=\alpha f(P_{t})=\alpha\,\psi(t).
\]
La funzione $\chi(t)=t^{-\alpha}\psi(t)$ ha allora, per ogni $t>0$, derivata
$\chi'(t)=t^{-\alpha-1}\bigl[t\,\psi'(t)-\alpha\,\psi(t)\bigr]=0$. Per il
\cref{cor:funzioni-costanti}, applicato nell'intervallo di estremi $1$ e $t$, risulta
$\chi(t)=\chi(1)=f(x,y)$, cioè $f(tx,ty)=t^{\alpha}f(x,y)$.
\end{dimostrazione}

\begin{esempio}[Sensibilità relative e grado di omogeneità]\label{es:omogenee-sensibilita}
Se $f>0$, dividendo la \eqref{eq:eulero-omogenee} per $f$ si ottiene
\[
  \frac{x\,f_{x}}{f}+\frac{y\,f_{y}}{f}=\alpha :
\]
\emph{le sensibilità relative di una funzione omogenea hanno per somma il grado}. Per il
coefficiente di resistenza indotta $C_{Di}=\frac{C_{L}^{2}}{\pi e\,AR}$, omogeneo di grado
$1$, le derivate calcolate nell'\cref{es:sensibilita-AR} danno
\[
  \frac{C_{L}}{C_{Di}}\,\frac{\partial C_{Di}}{\partial C_{L}}=2, \qquad
  \frac{AR}{C_{Di}}\,\frac{\partial C_{Di}}{\partial AR}=-1,
\]
la cui somma è proprio $1$. Con il differenziale,
$\frac{\Delta C_{Di}}{C_{Di}}\approx2\,\frac{\Delta C_{L}}{C_{L}}-\frac{\Delta AR}{AR}$: per volare
con un $C_{L}$ maggiore del $10\%$ senza aumentare il coefficiente di resistenza indotta
occorre un allungamento maggiore di circa il $20\%$. Allo stesso modo la portanza
$L(\rho,V,C_{L})$, omogenea di grado $1+2+1=4$, ha le sensibilità relative $1$, $2$, $1$
trovate nell'\cref{es:portanza-tre}: il teorema si estende infatti alle funzioni di $n$
variabili, nella forma $\sum_{i}x_{i}f_{x_{i}}=\alpha f$.
\end{esempio}

\needspace{8\baselineskip}
\begin{python}[Il teorema di Eulero verificato con sympy]
import sympy as sp

x, y = sp.symbols("x y", positive=True)
CL, AR, e = sp.symbols("C_L AR e", positive=True)
grado = lambda f, u, v: sp.simplify(       # per Eulero è il grado
    (u * sp.diff(f, u) + v * sp.diff(f, v)) / f)
for f in (x**2 + 3*x*y, sp.sqrt(x**2 + y**2), x*y / (x**2 + y**2),
          x / (x**2 + y**2)):
    print(f"f = {str(f):18}  (x f_x + y f_y)/f = {grado(f, x, y)}")
CDi = CL**2 / (sp.pi * e * AR)             # resistenza indotta
print("C_Di = C_L^2/(pi e AR): grado =", grado(CDi, CL, AR))
\end{python}
\begin{risultato}
f = x**2 + 3*x*y        (x f_x + y f_y)/f = 2
f = sqrt(x**2 + y**2)   (x f_x + y f_y)/f = 1
f = x*y/(x**2 + y**2)   (x f_x + y f_y)/f = 0
f = x/(x**2 + y**2)     (x f_x + y f_y)/f = -1
C_Di = C_L^2/(pi e AR): grado = 1
\end{risultato}

Le funzioni omogenee sono anche alla base dell'\emph{analisi dimensionale}: una legge
fisica non dipende dalle unità di misura, ed è quindi omogenea rispetto a ciascuna
grandezza fondamentale; così si ricava, per un'ala di forma e assetto assegnati e
trascurando viscosità e compressibilità, che la portanza è proporzionale a $\rho V^{2}S$.
